\documentclass{report}
\usepackage{amsmath,amsthm}
\newtheorem{theorem}{Teorema}[chapter]
\newtheorem{definition}[theorem]{Definizione}
\theoremstyle{definition}
\newtheorem{example}[theorem]{Esempio}
\title{Mini}
\begin{document}
\chapter{Variabili aleatorie continue}
\section{Densità}
Una variabile $X$ è \emph{continua} se esiste $f\ge 0$.\footnote{Per questo gli estremi non contano.}
\begin{definition}
Sia $f$ una densità.
\end{definition}
\begin{theorem}[Disuguaglianza di Čebyšëv]\label{thm:c}
Per ogni $k>0$
\begin{equation}
P(|X-\mu|\ge k\sigma)\le \frac1{k^2}.
\end{equation}
\end{theorem}
\begin{proof}
Si spezza l'integrale. Vedi il Teorema~\ref{thm:c}.
\end{proof}
\begin{example}
Esponenziale.
\end{example}
\section{Allineamenti}
\begin{align*}
a &= b\\ c &= d
\end{align*}
\chapter{Secondo}
Testo con riferimento al capitolo~\ref{thm:c}.
\end{document}
